\documentclass[10pt]{article}
\usepackage{titling}

\title{Acids and Bases}
\date{DATE PLACEHOLDER}
\author{Ingyu Bahng}

\usepackage[letterpaper, margin=1in]{geometry}
\usepackage{amsmath} % better math functions
\usepackage{amssymb} % better math symbols
\usepackage{babel}[english] % change rules to english
\usepackage{lipsum} % allows for lorem ipsum text generation
\usepackage{xr-hyper} % allows cross referencing labels from external LaTeX documents 
\usepackage{hyperref} % for hyperlinking text
\usepackage{fancyvrb} % for better verbatim environments
\usepackage{newverbs} % allows for inline typewriter font via \texttt{}
\usepackage{fancyvrb} % also code font blocks but has more capabilities than texttt
\usepackage{footnote} % allows footnotes inside tabulars
\usepackage{dcolumn} % allows for decimal alignment in tables
\usepackage{tabularx}

\usepackage{fontspec}
\setmainfont{Gentium Plus} % allowing greek letters in text

\usepackage{unicode-math}
\setmathfont{XITS Math} % allowing greek letters in math



% tcolorbox package
\usepackage{tcolorbox}
\tcbuselibrary{breakable}

\tcbset{ % this is the default design of tcolorboxes
  colframe = black,
  colback  = black!1,
  coltitle = black,
  colbacktitle = black!15,
  breakable, 
  arc=0mm,
  boxrule=0.5pt,
  toptitle=3pt,
  bottomtitle=3pt,
  left=8pt,
  right=8pt,
  top=6pt,
  bottom=6pt,
  before skip=12pt,
  after skip=12pt,
  bottomrule at break=-1pt,
  toprule at break=-1pt,
  fonttitle=\bfseries,
}

% Custom Environments
\newtcolorbox[auto counter, number within=section]{definition}[1][]
{
    title = \textbf{Definition \thetcbcounter: ~#1}
}

\newtcolorbox[auto counter, number within=section]{core}[1][]
{
    title = \textbf{Core Concept \thetcbcounter: ~#1}
}


\usepackage{fancyhdr}
\usepackage[skip=0.8em, indent=0em]{parskip}

\pagestyle{fancy} % this section generates the lines and text at the top and bottom of each page
\fancyhead[L]{\thetitle}
\fancyhead[C]{\theauthor}
\fancyhead[R]{\thedate}
\fancyfoot[C]{\thepage}
\renewcommand{\footrulewidth}{0.3pt}
\renewcommand{\thispagestyle}[1]{}

\renewcommand{\arraystretch}{1.5} % table vertical padding
\setlength{\tabcolsep}{10pt} % table horizontal padding

\usepackage{enumitem} % better control over enumerate and itemize

% List customization
\setlist[itemize]{%
  label=\bullet,
  itemsep=2pt,
  leftmargin=1.5em,
}

\setlist[enumerate]{%
  label=(\alph*),
  itemsep=2pt,
  leftmargin=1.5em,
}



\usepackage{pgfplots} % for plot creation
\pgfplotsset{compat=1.18} % setting the version used for pgfplots
\usepackage{float} % for better figure placement using [h]
\usepackage{graphicx} % for importing figures into the tex file
\usepackage{subcaption} % allows for multiple subfigures with individual captions wihtin one figure
\usepackage{xparse} % allows for more than 9 parameters in commands
\usepackage{adjustbox} % for valign option
\captionsetup{font=small} % setting the caption fontsize to small always

\usepackage{silence}
\WarningsOff[float]  % suppress float package warnings

\newcommand{\myfig}[4]{%
  \begin{figure}[H]
    \centering
    \adjustbox{valign=c}{\includegraphics[width=#1\textwidth]{#2}}
    \caption{#3}
    \label{#4}
  \end{figure}
}

\newcommand{\mymultifig}[8]{
  \begin{figure}[H]%
    \centering%
    \begin{subfigure}{#1\textwidth}%
      \centering%
      \includegraphics[width=\linewidth]{#2}%
    \end{subfigure}%
    \hfill%
    \begin{subfigure}{#3\textwidth}%
      \centering%
      \includegraphics[width=\linewidth]{#4}%
    \end{subfigure}%
    \hfill%
    \begin{subfigure}{#5\textwidth}%
      \centering%
      \includegraphics[width=\linewidth]{#6}%
    \end{subfigure}%
    \caption{#7}%
    \label{#8}%
  \end{figure}%
}

\usepackage{newverbs} % allows for inline typewriter font via \texttt{}
\usepackage{fancyvrb} % also code font blocks but has more capabilities than texttt
\usepackage{listings} % allows for codeblocks - config below

%BELOW ARE CODE BLOCK DESIGNS (Light Theme)
\definecolor{gray}{HTML}{707070}
\definecolor{lightgray}{HTML}{333333}
\definecolor{codebg}{HTML}{F5F5F5}
\definecolor{keyword}{HTML}{0000FF}
\definecolor{string}{HTML}{A31515}
\definecolor{number}{HTML}{098658}
\definecolor{function}{HTML}{795E26}
\definecolor{comment}{HTML}{008000}
\definecolor{classname}{HTML}{267F99}

% default options for listings (for code)
\lstset{
  frame=ltbr,
  language=python,
  aboveskip=3mm,
  belowskip=3mm,
  showstringspaces=false,
  columns=fullflexible,
  keepspaces=true,
  basicstyle=\fontsize{9}{10}\ttfamily\color{lightgray},
  numbers=left,
  firstnumber=1,
  numberstyle=\fontsize{7}{10}\ttfamily\color{gray},
  keywordstyle=\color{keyword},
  commentstyle=\color{comment},
  stringstyle=\color{string},
  backgroundcolor=\color{codebg},
  breaklines=true,
  breakatwhitespace=true,
  tabsize=2,
  xleftmargin=3em,
  numbersep=1em,
  framexleftmargin=2.5em,
  framesep=6pt,
  stepnumber=1,
}

\hfuzz=5.002pt % ignore overfull hbox badness warnings below this limit



\usepackage{chemfig}
\usepackage[version=4]{mhchem}

\newcommand{\bce}[1]{
  \hspace{7mm}{#1}
}

\newcommand{\sno}[0]{
  $\mathrm{S}_{\mathrm{N}}1$
}

\newcommand{\snt}[0]{
  $\mathrm{S}_{\mathrm{N}}2$
}

\newcommand{\reactionfig}[4]{
    \begin{figure}[H]
    \centering
    \scalebox{#4}{
    \schemestart
    #1
    \schemestop
    }
    \caption{#2}
    \label{#3}
    \end{figure}
}




\begin{document}

\input{../../presets/_general/startup}

\section{Acid-Base Basics}

  \begin{definition}[Acids]

    \textbf{Acids} are substances that increase the [\ce{H+}] of a solution.

    \bce{(shorthand) \ce{HA(aq) <=> H^{+}(aq) + A^{-}(aq)}}
    
    \bce{(more accurately) \ce{HA(aq) + \ce{H2O(l)} <=> H3O^{+}(aq) + A^{-}(aq)}}

  \end{definition}

  \begin{definition}[Bases]

    \textbf{Bases} are substances that decrease the [\ce{H+}] of a solution (usually through the addition of \ce{OH-}, however there are exceptions).

    \bce{(shorthand) \ce{BOH(aq) + \ce{H+}(aq) <=> B^{+}(aq) + H2O(aq)}}
    
    \bce{(more accurately) \ce{BOH(aq) + \ce{H3O+}(aq) <=> B^{+}(aq) + 2H2O(l)}} 
    
  \end{definition}
 
  \begin{core}[Brønsted-Lowry Theory]

     \textbf{Brønsted-Lowry Theory} defines acids as \ce{H+} donors and bases as \ce{H+} acceptors (not just \ce{OH-} donors).
     
  \end{core}

  As covered in \href{../equilibrium/equilibrium.pdf}{Equilibrium}, acid-base reactions also have equilibriums (\textit{denoted as $K_{a/b}$}) specific to each acid/base.

  \begin{definition}[Strong/Weak Acids and Bases]
   
    We categorize acid-base reactions with high $K_{a/b}$ as \textbf{strong acids/bases}
    \[
      K_{a/b} >> 1
    \]
    and those with low $K_{a/b}$ as \textbf{weak acids/bases}.
    \[
      K_{a/b} << 1
    \]
    Strong acids and bases also tend to conduct electricity better because of more dissociated ions.

  \end{definition}
  
  \begin{definition}[Indicators]

    \textbf{Indicators} (In) are weak acids or bases that change color at some [\ce{H+}] in a solution. The color change occurs because the cooresponding acid (\ce{HIn}) and base (\ce{In-}) forms of the indicators possess different colors.
  
    \bce{\ce{HIn(aq) <=> H^{+}(aq) + In^{-}(aq)}}
  
    thus $K_{In}$ is defined as 
    \[
      K_{In} = \frac{[H^{+}][In^{-}]}{[HIn]}
    \] 
    Here are a two examples of indicators that are commonly used in the lab setting.

    \bce{Anthocyanin: \ce{\textcolor{vRed}{HA}(aq) <=> H^{+}(aq) + \textcolor{vBlue}{$A^{-}$}(aq)}}

    \bce{Anthoxanthin: \ce{HB(aq) <=> H^{+}(aq) + \textcolor{vYellow}{$B^{-}$}(aq)}}

    Anthocyanin shows either a red or blue hue based on the comparative level of its acid form to its cooresponding base. Anthoxanthin, on the other hand, is either clear or yellow. Through Le Chatelier's principle, we can manipulate the color of the indicator solution by adding a strong acid or base to raise or lower [\ce{H^+}], respectively. 

  \end{definition}
  
\section{Acid-Base Stoichiometry}

  \begin{definition}[Amiphiproticity]
 
    Water is \textbf{amiphiprotic}, meaning that it can both accept or donate an \ce{H^+}.

    \bce{(\ce{H^+} Acceptor) \ce{HCl(aq) + \ce{H2O(l)} <=> H3O^{+}(aq) + Cl^{-}(aq)}}

    \bce{(\ce{H^+} Donor) \ce{NH3(aq) + \ce{H2O}(l) <=> NH4^{+}(aq) + OH^{-}(aq)}}

    Water also maintains equilibrium with the autoionized form of itself. Heat is a reactant in this autoionization reaction, thus an increase or decrease in temperature will shift the reaction right and left, respectively.

    \bce{\ce{H2O(l) + H2O(l) + \text{heat} <=> H3O^{+}(aq) + OH^{-}(aq)}}

    Therefore the equlibrium of \ce{H2O(l)} is represented as
    \[
      K_{W} = [H3O^{+}][OH^{-}] = [H^{+}][OH^{-}]
    \]

  \end{definition}

  \begin{definition}[pH]

    \textbf{pH} is a quantitative measure of how acidic (low pH) or basic (high pH) a solution is. Mathematically, pH is defined as 
    \[
      pH = -\log_{10}[H^{+}]
    \]  

  \end{definition}

  \begin{table}[H]
  \centering
  \begin{tabular}{lccc}
  \hline
  Water temperature (\textdegree C) & $K_w$ & $[\text{H}^+]$ & pH \\
  \hline
  0  & $10^{-14.94}$ & $10^{-7.47}$ & 7.47 \\
  5  & $10^{-14.73}$ & $10^{-7.36}$ & 7.36 \\
  10 & $10^{-14.53}$ & $10^{-7.26}$ & 7.26 \\
  15 & $10^{-14.35}$ & $10^{-7.18}$ & 7.18 \\
  20 & $10^{-14.17}$ & $10^{-7.08}$ & 7.08 \\
  25 & $10^{-14.00}$ & $10^{-7.00}$ & 7.00 \\
  30 & $10^{-13.83}$ & $10^{-6.92}$ & 6.92 \\
  35 & $10^{-13.68}$ & $10^{-6.84}$ & 6.84 \\
  40 & $10^{-13.53}$ & $10^{-6.76}$ & 6.76 \\
  \hline
  \end{tabular}
  \caption{Water pH across different temperatures}
  \label{tab:water_ph}
  \end{table} 

  A common misconception people often assume is that the pH of 7 corresponds to a neutral solution of water (meaning $[\ce{H+}] = [\ce{OH-}]$). From \autoref{tab:water_ph} however, we see that \textbf{the pH of 7 being neutral only applies to a T of 25°C}. At higher temperatures, the Eq will shift right to produce more \ce{H3O+}, effectively making a lower pH value as the new neutral base point.

  From a molecular perspective, higher temperatures increases autoionization rates because the higher KE of the water molecules makes it statistically more likely to collide/react with one another in order to form \ce{H3O+} and \ce{OH-}, thus decreasing the pH of the solution.

  Let's take a look at stoichiometric calculations for [\ce{H^+}] for strong acids and bases.

  For \textbf{strong acids and bases}, we rely on simple stoichiometry to calculate [\ce{H^+}]. We'll use HCl as an example.

  \bce{\ce{HCl(aq) + \ce{H2O(l)} <=> H3O^{+}(aq) + Cl^{-}(aq)}}

  The equilibrium for HCl is represented as 
  \[
    K_{a} = \frac{[\ce{H3O+}][Cl^{-}]}{[HCl]} >> 1
  \]

  The $K_{a/b}$ for strong acids and bases are generally very large values. The approximate $K_a$ for HCl is $10^{7}$. Thus when dealing with strong acids and bases, we can assume the initial [HCl] to be equal to the final [$H3O^{+}$] with a negligible margin of error.
  \[
    [HCl] = [\ce{H3O+}] = [\ce{H^+}]
  \]

  Subsequently, you can calculate pH using the formala $pH = -\log_{10}[H^{+}]$

  For \textbf{weak acids and bases}, we rely on more complex stoichiometry to calculate [\ce{H^+}]. We'll use acetic acid (\ce{CH3COOH}) as an example.

  \bce{\ce{CH3COOH(aq) + \ce{H2O(l)} <=> H3O^{+}(aq) + CH3COO^{-}(aq)}}

  The equilibrium for acetic acid is represented as 
  \[
    K_{a} = \frac{[H3O^{+}][\ce{CH3COO-}]}{[\ce{CH3COOH}]} << 1
  \]

  The $K_{a/b}$ for weak acids and bases are generally very small values. The approximate $K_a$ for acetic acids is $1.8 \times 10^{-5}$. Thus when representing the $K_a$ of acetic acid, we can derive the equation
  \[
    K_a = 1.8 \times 10^{-5} = \frac{[H3O^{+}]_{final}[\ce{CH3COO-}]_{final}}{[\ce{CH3COOH}]_{initial} - [H3O^{+}]_{final}} \approx \frac{[H3O^{+}]_{final}[\ce{CH3COO-}]_{final}}{[\ce{CH3COOH}]_{initial}}
  \]

  Here, since a mole of \ce{CH3COOH} is consumed every mole of \ce{H3O+} or \ce{CH3COO-} is formed, $[\ce{CH3COOH}]_{final}$ must be equal to $[\ce{CH3COOH}]_{initial} - [H3O^{+}]_{final}$. However, since $K_a$ is extremely small, we approximate $[\ce{CH3COOH}]_{final}$ as just $[\ce{CH3COOH}]_{initial}$ because $[H3O^{+}]_{final}$ value is negligible in comparison to $[\ce{CH3COOH}]_{initial}$.

  The smaller the $K_{a/b}$ of a weak acid or base is, the more accurate the stoichiometric calculation will be.

  \begin{definition}[Conjugate Acid-Base Pairs]

    The cooresponding base that is formed when an acid donates a \ce{H^+} is called the \textbf{conjugate base}. Likewise, the cooresponding acid that forms when a base accepts a \ce{H^+} is called the \textbf{conjugate acid}.

    For example, the conjugate acid for \ce{NH3} is \ce{NH4+}, and the conjugate base for \ce{HNO3} is \ce{NO3-}. These are represented as

    \bce{\ce{NH3(aq) + H+(aq) <=> NH4^{+}(aq)}}

    \bce{\ce{HNO3(aq) <=> H+(aq) + NO3^{-}(aq)}}

    Because the chemical reaction for an acid is the same as that of the conjugate base (\textit{just reversed}), the $K_a$ for an acid is the inverse of the $K_b$ for its conjugate base. As such, the stronger the acid/base, the weaker the corresponding conjugate base/acid will be.  
    
  \end{definition}

  \begin{core}[Conjugate Acid-Base Pair Equilibriums and $K_w$]

    The product of the equilibrium constants of an acid $K_a$ and its conjugate base $K_b$ (\textit{or base and its conjugate acid}) is equal to the equilibrium constant of the autoionization of water $K_w$. We can derive this relationship with the following algebraic manipulations using the formulas for the shell reactions below.

    \bce{(Acid) \ce{HA(aq) + H2O(l) <=> H3O+(aq) + A-(aq)}}

    \bce{(Conjugate Base) \ce{A-(aq) + H2O(l) <=> HA(aq) + OH-(aq)}}  

    Thus the equilibrium constants are $K_{a} = \frac{[H3O^{+}][\ce{A-}]}{[\ce{HA}]} , K_{b} = \frac{[\ce{HA}][OH^{-}]}{[\ce{A-}]}$. When we rewrite $K_b$ as
    \[
      K_{b} = \frac{[\ce{HA}][OH^{-}]}{[\ce{A-}]} \times \frac{[\ce{H3O+}]}{[\ce{H3O+}]} = K_w/K_a \Rightarrow K_w = K_a K_b
    \]
    This relationship is true for all conjugate acid-base pairs. 
     
  \end{core}

\section{Titration and Buffers}

  \begin{definition}[Titration]

    \textbf{Titration} is the processs of adding an acid or base to bring it about to a level where occurs a change at a previously known pH (via indicators). This laboratory technique \textbf{allows chemists to determine the pH of a solution} by observing the amount of base or acid it takes to reach the target pH level.
   
    \myfig{0.55}{acids_and_bases_figures/titration.png}{Basic titration set up}{fig:titration}

    We call the solution that is being titrated the \textbf{analyte} and the solution that is doing the titrating the \textbf{titrant}.

  \end{definition}

  \begin{core}[Strong Acid and Strong Base Titration]

    An example of a \textbf{strong acid} (HCl) being titrated with a \textbf{strong base} (NaOH) is represented by the chemical equation \ce{HCl(aq) + NaOH(aq) <=> NaCl(aq) + H2O(l)}. Because of their high $K_{a/b}$ values, we can simplify this equation to the following.

    \bce{\ce{H+(aq) + OH-(aq) <=> H2O(l)}}

    To neutralize a solution of an unknown [HCl], we would need to add equal amounts of NaOH. Thus, based on the amount of NaOH we add to the solution in order to reach this neutral point, we can calculate the original [HCl].

    \myfig{0.4}{acids_and_bases_figures/sa_sb_tcurve.png}{Strong acid and base titration curve}{fig:sa_sb_tcurve}

    \autoref{fig:sa_sb_tcurve} is a graphical representation of the analyte's pH as more titrant is added. The point at which the moles of titrant added is stoichiometrically equal to the amount of the analyte originally present is called the \textbf{equivalence point}. For the case of a strong acid with a strong base, this point is at a neutral pH of 7 at 25$^\circ$C.

  \end{core}

  \begin{core}[Weak Acid/Base and Strong Base/Acid Titration]

    The titration curve of a \textbf{weak acid} (\ce{CH3COOH}) being titrated by a \textbf{strong base} (NaOH) is shown below.

    \myfig{0.4}{acids_and_bases_figures/wa_sb_tcurve.png}{Weak acid and strong base titration curve}{fig:wa_sb_tcurve}

    The equivalence point for this case \textbf{lies above the neutral pH}, indicating a basic state at equal parts weak acid and strong base. This is because once all \ce{H^+} is neutralized by the presence of \ce{OH-} from NaOH, \textbf{what is left is the strong conjugate base molecule} \ce{CH3COO-}. In the presence of this strong conjugate base, some water molecules will react with the \ce{CH3COO-} to form \ce{CH3COOH} and \ce{OH-}, thus resulting in a higher pH value at equivalance. When confused, it's helpful to (1) assume the reaction will go to completion and then (2) look at what is left and ask how it affects the pH.

  \end{core}

  \begin{definition}[Buffers]
     
    \textbf{Buffers} are solutions containing both a weak acid and its cooresponding conjugate base. They are used to \textbf{resist} large changes of pH upon addition of strong acids or bases.

    Based on the chemical equation of a weak acid (\ce{HA(aq) + H2O(l) <=> A-(aq) + H3O+(aq)}), we can manipulate the $K_a$ equation to give us the final following formula.
    \[
      K_a = \frac{[H3O^{+}][\ce{A-}]}{[\ce{HA}]} = [H3O^{+}]\frac{[\text{conjugate base}]}{[\text{conjugate acid}]}
    \]
    \[
      -\log K_a = -\log [H3O^{+}] + \log \frac{[\text{conjugate base}]}{[\text{conjugate acid}]}
    \]
    \[
      pH = pK_a + \log \frac{[\text{conjugate base}]}{[\text{conjugate acid}]}
    \]
    The final equation is a simple way to determine the pH of a buffer solution under a given concentration of the weak acid and its conjugate base. \textit{Notice how the pH of a solution will be equivalent to the $pK_a$ when there are equal parts conjugate acid and base}

  \end{definition}

  The making a buffer solution can be done in one of three ways.

  \begin{enumerate}
      \item Equal parts of the weak acid and its conjugate base
      \item Weak acid with half the amount of a strong base
      \item Conjugate base with half the amount of a strong acid 
  \end{enumerate}

  \myfig{0.35}{acids_and_bases_figures/buffer_curve.png}{Buffer solution titration curve}{fig:buffer_curve}

  Based on \autoref{fig:buffer_curve} the $pK_a$ is also half way between the equivalance point and y-axis of the titration curve, the point at which $[\text{conjugate acid}] = [\text{conjugate base}]$. Any subsequent addition of a strong acid or base will shift the pH value, as shown in the following titration curve, but the pH will be resistant to change within a given range until larger amounts of a strong acid or base is added.

  Using this technique, chemists can customize the pH they want a solution to linger around based on the $K_a$ value of a weak acid. Conversely, a buffering solution can also be held at a basic pH level by adding equal parts of weak base with its cooresponding strong conjugate acid.

  
\end{document}
